\subsection{Suma de Prefijos con Mapa de Primera Aparición}

\graybold{Preguntas guía:}

\begin{itemize}
\item ¿Piden el subarreglo contiguo más largo que cumple una condición de balance (igual cantidad de dos tipos de elementos o suma igual a un valor fijo)?
\item ¿Puedo transformar cada elemento en un valor (+1, -1 o 0) para que la condición se convierta en que la suma del subarreglo sea igual a 0?
\item ¿Basta con almacenar únicamente la primera posición donde apareció cada suma de prefijo?
\end{itemize}

\graybold{Utilidad:}

Encontrar en tiempo lineal el subarreglo contiguo más largo cuya suma (o balance entre dos tipos de elementos) sea igual a un valor fijo, evitando revisar todos los \pow{n}{2} subarreglos posibles. Es una técnica muy utilizada en problemas de balances, diferencias entre cantidades y prefijos acumulados.

\graybold{Complejidad:}

Tiempo \bigO{n}. Espacio \bigO{n}.

\graybold{Fundamento teórico:}

Sea $P[i]$ la suma de prefijos de los primeros $i$ elementos del arreglo. La suma del subarreglo $(i+1 \dots j)$ puede expresarse como
\[
P[j]-P[i].
\]

Por lo tanto, un subarreglo tiene suma igual a cero si y solo si
\[
P[j]=P[i].
\]

Para obtener el subarreglo de mayor longitud basta con almacenar, para cada valor de suma acumulada, únicamente la \textbf{primera} posición donde apareció. Cada vez que el mismo prefijo vuelve a aparecer, el segmento comprendido entre ambas posiciones cumple la condición buscada.

\graybold{Ejercicio aplicado}

Dada una secuencia de $N$ transacciones (positiva = depósito, negativa = retiro y cero = consulta), calcular:

\begin{itemize}
\item El depósito de mayor valor.
\item El retiro más negativo.
\item La longitud del subperíodo contiguo más largo que contiene igual cantidad de depósitos y retiros.
\end{itemize}

\graybold{I/O:}

$N \le 10\,000$.

Existen múltiples casos de prueba terminados por el centinela \texttt{0}. Debido a que cada caso consiste en leer un entero seguido de un arreglo, la lectura mediante \texttt{sys.stdin.buffer.readline()} resulta más eficiente que leer toda la entrada de una vez.

\graybold{Código solución}

\begin{lstlisting}[language=Python]
import sys

input = sys.stdin.buffer.readline

def solve():
    out = []
    while True:
        n = int(input())

        if n == 0:
            break
        arr = list(map(int, input().split()))

        deposit = 0
        withdrawal = 0

        for x in arr:
            if x > 0:
                deposit = max(deposit, x)
            elif x < 0:
                withdrawal = min(withdrawal, x)

        first_seen = {0: 0}
        prefix = 0
        longest = 0

        for i, x in enumerate(arr, start=1):
            if x > 0:
                prefix += 1
            elif x < 0:
                prefix -= 1

            if prefix in first_seen:
                longest = max(longest, i - first_seen[prefix])
            else:
                first_seen[prefix] = i

        out.append(f"{deposit} {withdrawal} {longest}")
    sys.stdout.write("\n".join(out) + "\n")
solve()
\end{lstlisting}